\documentclass[11pt]{article}

\usepackage[a4paper,margin=2.5cm]{geometry}
\usepackage[utf8]{inputenc}
\usepackage{xcolor}
\usepackage{listings}
\usepackage{hyperref}
\usepackage{booktabs}
\usepackage{parskip}
\usepackage{enumitem}
\usepackage{titlesec}
\usepackage{microtype}
\usepackage[htt]{hyphenat}

% Long inline \texttt{} paths/identifiers (file paths, CMake flags) have no
% natural break points and were overflowing past the right margin. microtype
% + hyphenat[htt] lets LaTeX hyphenate/break inside \texttt text, and
% emergencystretch is a last-resort fallback so a genuinely unbreakable token
% loosens inter-word spacing instead of overflowing the page.
\emergencystretch=3em
\sloppy

\hypersetup{
  colorlinks=true,
  linkcolor=blue!50!black,
  urlcolor=blue!50!black,
  citecolor=blue!50!black,
}

\definecolor{codebg}{RGB}{245,245,245}
\lstset{
  basicstyle=\ttfamily\small,
  backgroundcolor=\color{codebg},
  columns=fullflexible,
  breaklines=true,
  frame=single,
  framesep=4pt,
  xleftmargin=2pt,
  xrightmargin=2pt,
  showstringspaces=false,
}

\titleformat{\section}{\normalfont\Large\bfseries}{\thesection}{1em}{}
\titleformat{\subsection}{\normalfont\large\bfseries}{\thesubsection}{1em}{}

\title{GEM-garfield-simulation \\ \large Installation Guide}
\author{}
\date{Last verified: 2026-09-26}

\begin{document}
\maketitle

\tableofcontents

\section{Overview}

This document is a detailed, step-by-step installation guide for the
software stack that \texttt{GEM-garfield-simulation} depends on:
\textbf{Gmsh} (geometry/mesh generation), \textbf{Elmer FEM}
(electrostatic field solve), \textbf{Garfield++/Magboltz} (microscopic
electron avalanche simulation), \textbf{ROOT} (I/O and, indirectly, the
ABI that Garfield++ must be built against), and the \textbf{Python}
environment used for analysis and visualization.

\texttt{README.md} in the repository root gives the minimum a returning
user needs (a version table and where paths are set); this document is
the fuller reference for someone setting this stack up for the first
time, or rebuilding a piece of it after an upgrade.

This is written for a shared HEP cluster (e.g.\ KEKCC), where you
typically have no sudo/admin rights. Every install step below goes under
\texttt{\$HOME}, not a system directory.

\subsection{Why installation order matters}

The one constraint that drives most of the pain in this stack is:
\textbf{Garfield++ must be built against the exact same ROOT installation
(same version, same compiler ABI) that everything else in the pipeline
uses.} ROOT's dictionary-generation ABI changes between versions, so a
Garfield++ build linked against ROOT~6.32 will fail to link against a
ROOT~6.40 executable with errors like:

\begin{lstlisting}
undefined reference to `ROOT::TGenericClassInfo::TGenericClassInfo(...)'
undefined reference to `TGeoBBox::ComputeNormal(...)'
\end{lstlisting}

There is no way around this other than keeping one consistent ROOT
version and rebuilding Garfield++ against it whenever that version
changes. The recommended order is therefore:

\begin{enumerate}
  \item Install ROOT first, and settle on the version you intend to keep.
  \item Install Gmsh and Elmer (independent of ROOT/Garfield++).
  \item Build Garfield++ against that same ROOT installation.
  \item Set up the Python environment (independent, but see
    Section~\ref{sec:python} for a real package-conflict gotcha).
  \item Build this repository's own C++ macros against the same
    ROOT + Garfield++ pair.
\end{enumerate}

\subsection{Versions verified to work together}

The following combination is what this project is actually built and
tested against (2026-09-26). Later point releases of any one component
are not guaranteed to be compatible with the others without re-testing.

\begin{table}[h]
\centering
\begin{tabular}{@{}p{6.5cm}p{5cm}@{}}
\toprule
Tool & Version \\
\midrule
Gmsh (CLI) & 4.13.1 \\
Gmsh Python API (\texttt{pip install gmsh}) & 4.15.2 \\
Elmer FEM & v9.0 \\
Garfield++ & package version string ``0.3'' (built against ROOT below) \\
ROOT & 6.40.04 \\
CMake & 3.31.8 \\
gcc & 11.5.0 (Red Hat 11.5.0-14) \\
\bottomrule
\end{tabular}
\end{table}

\subsection{System-level compiler prerequisites}

Before building anything below, confirm the system has:
\begin{itemize}
  \item a C++ compiler supporting C++17 (\texttt{gcc}/\texttt{g++} 11.5.0
    here; any reasonably recent gcc/clang works),
  \item a Fortran compiler (\texttt{gfortran}), and
  \item CMake $\geq$ 3.16.
\end{itemize}
The Fortran compiler is easy to miss: Garfield++'s own
\texttt{CMakeLists.txt} declares \texttt{project(... LANGUAGES Fortran
CXX)} unconditionally, since Magboltz (the gas-transport code Garfield++
calls into) is written in Fortran -- the Garfield++ build fails
immediately at the \texttt{cmake} step without one, with an error
naming the Fortran compiler check.

\textbf{GSL (GNU Scientific Library) is not needed here}, despite some
older Garfield++ documentation/notes listing it as a prerequisite --
confirmed by reading Garfield++'s own \texttt{CMakeLists.txt}
(\texttt{option(GARFIELD\_WITH\_GSL "Build Garfield++ with GSL (useless
if C++>=17)" OFF)}, only switched on and required when building with a
C++ standard older than 17). This project's Garfield++ build already
uses C++17 throughout, so GSL is neither installed nor linked.

\section{ROOT}

Most HEP clusters (KEKCC and similar) already have a ROOT installation
available. Use that one -- the only project-specific requirement is that
Garfield++ (Section~5) gets built against the exact same ROOT. Check the
version with \texttt{root-config --version}.

If no ROOT install is available, build it yourself, no sudo/admin rights
needed (everything goes under \texttt{\$HOME}):

\begin{lstlisting}
git clone --branch v6-40-04 --depth 1 https://github.com/root-project/root.git root_src
mkdir root_build && cd root_build
cmake -DCMAKE_INSTALL_PREFIX=$HOME/local/root/6.40.04 ../root_src
cmake --build . -j"$(nproc)"
cmake --install .
\end{lstlisting}

Full build instructions/options:
\url{https://root.cern/install/build_from_source/}.

This project does not rely on \texttt{thisroot.sh} being sourced
globally (e.g. from \texttt{.bashrc}) -- each script/build step sets the
paths it needs explicitly (see Section~\ref{sec:paths}). Just be
consistent about which ROOT installation every other tool below is
pointed at.

\section{Gmsh}

Two independent things are needed: the Gmsh CLI (only used incidentally;
this project drives Gmsh almost entirely through its Python API) and the
Gmsh Python module.

\subsection{Gmsh CLI (optional)}

Download a prebuilt binary from \url{https://gmsh.info/#Download} and
extract it, e.g.:

\begin{lstlisting}
mkdir -p ~/local
cd ~/local
wget https://gmsh.info/bin/Linux/gmsh-4.13.1-Linux64.tgz
tar xf gmsh-4.13.1-Linux64.tgz
~/local/gmsh-4.13.1-Linux64/bin/gmsh --version
\end{lstlisting}

\subsection{Gmsh Python API (required)}

\begin{lstlisting}
pip install gmsh
python3 -c "import gmsh; print(gmsh.__file__)"
\end{lstlisting}

The pip package's own bundled version (4.15.2 here) does not need to
exactly match the standalone CLI's version above; they are independent
installs and this project only actually calls the Python API at runtime.

\textbf{Import-order gotcha:} in some conda/pip environments, importing
\texttt{gmsh} before \texttt{matplotlib} in the same Python process
causes
\begin{lstlisting}
ImportError: /usr/lib64/libstdc++.so.6: version `CXXABI_1.3.15' not found
\end{lstlisting}
Gmsh's native extension can pull in the system's older
\texttt{libstdc++.so.6} first; once that specific shared object is loaded
process-wide, matplotlib's compiled extension (which needs a newer
libstdc++ ABI) can no longer load the version it actually needs. The fix
is purely at the Python-import level: \textbf{always \texttt{import
matplotlib} (or any module that imports it) before \texttt{import gmsh}}
in any script that uses both. This project's own scripts already follow
this convention; keep it if you add a new one that touches both.

\section{Elmer FEM}

Build from source (\url{http://www.elmerfem.org/blog/binaries/} also has
prebuilt binaries for some platforms, but a from-source build is more
portable across clusters):

\begin{lstlisting}
git clone https://github.com/ElmerCSC/elmerfem.git elmer_src
mkdir elmer_build && cd elmer_build
cmake -DCMAKE_INSTALL_PREFIX=$HOME/elmer \
      -DWITH_MPI=OFF \
      ../elmer_src
cmake --build . -j"$(nproc)"
cmake --install .
\end{lstlisting}

\texttt{-DWITH\_MPI=OFF} is a deliberate simplification here: this
project's meshes are solved as a single linear system per run (not
domain-decomposed), so MPI parallelism inside \texttt{ElmerSolver} is not
needed, and a non-MPI build avoids an entire class of shared-library
version-mismatch problems (e.g. a system \texttt{libmpi\_usempif08.so}
mismatch) that a shared/central Elmer install on a cluster can otherwise
run into. If your own environment already has a working MPI-enabled Elmer
install, that works too; just confirm \texttt{ElmerSolver} actually runs
(some shared HPC installs turn out to be broken in ways that only surface
at runtime, not at link/install time).

Verify with:
\begin{lstlisting}
export ELMER_HOME=$HOME/elmer
export PATH="$ELMER_HOME/bin:$PATH"
export LD_LIBRARY_PATH="$ELMER_HOME/lib:${LD_LIBRARY_PATH:-}"
ElmerSolver -version
\end{lstlisting}

\section{Garfield++}

Build from the official source
(\url{https://gitlab.cern.ch/garfield/garfieldpp}), pointed explicitly at
the ROOT installation from Section~2:

\begin{lstlisting}
git clone https://gitlab.cern.ch/garfield/garfieldpp.git garfieldpp_src
mkdir -p garfieldpp_src/build && cd garfieldpp_src/build

export ROOTSYS=$HOME/local/root/6.40.04
export PATH="$ROOTSYS/bin:$PATH"
export LD_LIBRARY_PATH="$ROOTSYS/lib:$LD_LIBRARY_PATH"

cmake -DROOT_DIR="$ROOTSYS/cmake" \
      -DCMAKE_INSTALL_PREFIX=$HOME/local/garfield \
      ..
cmake --build . -j"$(nproc)"
cmake --install .
\end{lstlisting}

\subsection{Rebuilding after a ROOT upgrade}

If ROOT is ever upgraded, Garfield++ \textbf{must} be rebuilt against the
new installation (rerun the same \texttt{cmake}/\texttt{build}/
\texttt{install} sequence above with \texttt{ROOTSYS} pointed at the new
ROOT) -- reusing the old \texttt{libGarfield.so} against a newer ROOT is
exactly the ABI mismatch described in Section~1.1. If the Garfield++
install is shared with other projects, confirm with anyone else relying
on it before rebuilding it in place.

\subsection{A CMake gotcha worth knowing about in advance}

On at least one system, a fresh \texttt{cmake ..} for a project
consuming Garfield++ (including this repository's own
\texttt{macros/CMakeLists.txt}) can silently resolve \texttt{ROOT\_DIR}
back to a stale/wrong ROOT installation, even after
\texttt{CMAKE\_PREFIX\_PATH} is set correctly. This happens because
Garfield's own \texttt{GarfieldConfig.cmake} internally calls
\texttt{find\_dependency(ROOT 6.0 COMPONENTS Geom Gdml)}, which can win
the resolution race depending on what \texttt{CMAKE\_PREFIX\_PATH} was
inherited from the environment. If \texttt{cmake ..}'s output reports the
wrong ROOT under a line like \texttt{Found Vdt: ...}, force it
explicitly:
\begin{lstlisting}
cmake -DROOT_DIR=$HOME/local/root/6.40.04/cmake ..
\end{lstlisting}
or add \texttt{set(ROOT\_DIR "..." CACHE PATH "" FORCE)} to the
consuming project's own \texttt{CMakeLists.txt} right after its
\texttt{CMAKE\_PREFIX\_PATH} is set (this repository's
\texttt{macros/CMakeLists.txt} already does this -- see that file's own
comments).

\section{Python environment}
\label{sec:python}

Any Python 3 environment works (conda, venv, or plain \texttt{pip install
--user}); this project does not depend on a specific environment manager.
No sudo needed either way -- a conda/venv environment installs under
\texttt{\$HOME} by default, and \texttt{--user} does the same for a
system Python:

\begin{lstlisting}
pip install --user gmsh numpy matplotlib uproot pyvista plotly
pip install --user "vtk==9.3.1"
\end{lstlisting}

\subsection{The \texttt{vtk} version pin (required)}

Installing \texttt{pyvista} alone pulls in whatever the latest
\texttt{vtk} wheel is. As of 2026-09, that is \texttt{9.7.x}, which is
\textbf{incompatible} with \texttt{trame\_vtk} (the library
\texttt{pyvista}'s HTML export path uses) and fails with:
\begin{lstlisting}
TypeError: unhashable type 'VTKAOSArray_vtkFloatArray'
\end{lstlisting}
while serializing a scene for HTML export (used by
\texttt{visualization/plot\_triple\_gem.py}). Explicitly pinning
\texttt{vtk==9.3.1} (after installing \texttt{pyvista}, so the pin takes
effect) resolves this. If a future \texttt{trame\_vtk} release fixes the
incompatibility, this pin can likely be relaxed -- but treat that as a
deliberate, tested decision, not an assumption.

\subsection{Versions this project is actually tested against}

\begin{table}[h]
\centering
\begin{tabular}{@{}ll@{}}
\toprule
Package & Version \\
\midrule
numpy & 2.4.4 \\
matplotlib & 3.10.9 \\
uproot & 5.7.4 \\
pyvista & 0.49.0 \\
plotly & 7.1.0 \\
vtk & 9.3.1 (pinned, see above) \\
gmsh (Python) & 4.15.2 \\
\bottomrule
\end{tabular}
\end{table}

Note: \texttt{uproot} (not PyROOT) is what this project's Python code
uses to read ROOT files. PyROOT is deliberately not used here, so the
Python environment's own version does not need to match ROOT's Python
binding at all -- one less version constraint to track.

\section{Building this project's C++ macros}

Once ROOT, Garfield++, and Elmer are installed and consistent with each
other:

\begin{lstlisting}
cd GEM-garfield-simulation/macros
mkdir -p build && cd build
cmake ..
cmake --build . -j"$(nproc)"
\end{lstlisting}

\texttt{macros/CMakeLists.txt} points \texttt{CMAKE\_PREFIX\_PATH} at
\texttt{\$HOME/local/garfield} and \texttt{\$HOME/local/root/6.40.04} by
default -- edit those two paths at the top of that file if your own
installation locations differ from this project's own development
machine.

If \texttt{cmake ..} reports the wrong ROOT (see
Section~5.2), pass \texttt{-DROOT\_DIR=...} explicitly as shown there.

\section{Where paths are set (no dependence on shell startup files)}
\label{sec:paths}

This project deliberately does \textbf{not} rely on \texttt{.bashrc} (or
similar) exporting \texttt{ROOTSYS}, \texttt{ELMER\_HOME}, or
\texttt{GARFIELD\_HOME} globally. Each entry point sets what it needs
explicitly:

\begin{itemize}
  \item \texttt{elmer/run\_field\_solve.sh} exports
    \texttt{ELMER\_HOME}/\texttt{PATH}/\texttt{LD\_LIBRARY\_PATH} itself
    before invoking \texttt{ElmerGrid}/\texttt{ElmerSolver}.
  \item \texttt{macros/CMakeLists.txt} sets \texttt{CMAKE\_PREFIX\_PATH}
    (and, if needed, \texttt{ROOT\_DIR}) itself at configure time.
  \item Python scripts read ROOT files via \texttt{uproot} (pure Python),
    so they need no ROOT environment at all.
\end{itemize}

If you set up your own equivalents of these paths differently from the
defaults baked into this repository, the two files above are the only
places that need editing -- there is no third, hidden place any of this
is duplicated.

\subsection{Optional: speeding up shell startup on a networked home directory}

If your home directory is on a networked filesystem (e.g. GPFS/NFS) and
you notice a large conda environment or ROOT installation making the
\emph{first} \texttt{python}/\texttt{root} invocation of a session slow
(tens of seconds, due to many small-file stat/open calls over the
network), packing the environment into a single squashfs image and
mounting it read-only under a node-local \texttt{/tmp} path is one
effective fix. This project's own development machine uses exactly this
approach (see \texttt{\textasciitilde/local/envfs\_README.md} outside
this repository for that machine's specific implementation) -- it is a
pure performance optimization, entirely optional, and orthogonal to
everything else in this guide: batch (LSF) jobs are unaffected by it
either way, since they always see the real, unmounted paths.

\section{Verifying the full stack end-to-end}

Once everything above is installed, confirm the whole pipeline actually
runs, from geometry generation through to Python analysis, using the
\texttt{triple\_gem\_field} model (this exact command sequence is the
one verified end-to-end in \texttt{docs/reference.md} \S2):

\begin{lstlisting}
cd GEM-garfield-simulation

cd geometry && python3 build_triple_gem_field_mesh.py && cd ..
cd elmer && bash run_field_solve.sh triple_gem_field && cd ..

cd macros/build
./export_avalanche_trajectories ../../results/mesh/triple_gem_field \
  ../../resources/ar_ch4_90_10.gas 5 -0.2029 0.8405 0.4235 0.021 \
  0.03637306695894642 0.1 0.0005 ../../results/root
cd ../..

python3 visualization/plot_triple_gem.py triple_gem_field
python3 geometry/analyze_plane_crossings.py \
  results/root/triple_gem_field_avalanche.root
\end{lstlisting}

A successful run produces
\texttt{results/root/triple\_gem\_field\_avalanche.root},
\texttt{results/html/triple\_gem\_field\_overview.html}, and prints a
funnel/summary from the last command with no Python traceback. See
\texttt{docs/reference.md} for what each pipeline stage's CLI arguments
mean (including how to compute them for a different mesh/model) and how
to run the full 7\texttimes{}7 production configuration.

\section{Summary of known compatibility pitfalls}

\begin{itemize}
  \item \textbf{Missing Fortran compiler}: Garfield++'s build fails
    immediately at \texttt{cmake} without \texttt{gfortran} (Magboltz is
    Fortran); GSL, despite older documentation, is not needed with a
    C++17 build (Section~1.3).
  \item \textbf{Garfield++ / ROOT ABI}: must be built against the exact
    same ROOT version; rebuild Garfield++ after any ROOT upgrade
    (Section~3.1).
  \item \textbf{CMake resolving the wrong ROOT\_DIR}: force it explicitly
    if a fresh \texttt{cmake ..} picks the wrong one (Section~5.2).
  \item \textbf{\texttt{gmsh} vs \texttt{matplotlib} import order}: always
    import matplotlib-touching modules before \texttt{gmsh} in the same
    process (Section~3, Gmsh).
  \item \textbf{\texttt{vtk} version}: pin to \texttt{9.3.1}; the latest
    wheel breaks \texttt{pyvista}'s HTML export (Section~6.1).
  \item \textbf{Elmer with MPI on a shared/HPC install}: confirm
    \texttt{ElmerSolver} actually runs at runtime, not just that it built
    -- a broken shared MPI runtime can pass CMake/link steps and still
    fail with a shared-library error the first time it is actually
    invoked (Section~4).
\end{itemize}

For gotchas encountered building and running the pipeline itself (as
opposed to installing the underlying tools), see
\texttt{docs/pipeline\_gotchas.md} in this repository.

\end{document}
